\documentclass{article}
\usepackage[utf8]{inputenc}
\usepackage{amsmath,amssymb}
\usepackage[most]{tcolorbox}
\usepackage{geometry}
\geometry{margin=2.5cm}

\newcommand{\step}[1]{\par\noindent\hangindent=3.5em\hangafter=1%
  \makebox[3.5em][l]{\textbf{#1.}}}
\newcommand{\steptag}[1]{\hfill{\small\textsf{[#1]}}}

\newtcolorbox{proofbox}[1]{
  colback=gray!5, colframe=gray!60,
  fonttitle=\bfseries, title={#1},
  breakable, enhanced,
  left=4pt, right=4pt, top=4pt, bottom=4pt,
  before skip=10pt, after skip=10pt
}

\begin{document}

\begin{proofbox}{There are universal e\_iso\textasciicircum{}r > 0, r\_iso > 0 such that, for...}
\step{1} There are universal e\_iso\textasciicircum{}r > 0, r\_iso > 0 such that, for every finite-dimensional exact-unit epsilon\_r-C*-algebra with 0 <= epsilon\_r <= e\_iso\textasciicircum{}r, J and -J are the only fixed points of the smooth sigma in their respective ambient r\_iso-balls. \steptag{validated} \\[6pt]
\step{1.1} Fix the universal tuple W supplied by lem-stage1-polar-constant-ledger and set e\_iso\textasciicircum{}r := C\_rect*e\_S1 and rho := r\_iso (the tuple field). These are positive universal constants. If 0 <= epsilon\_r <= e\_iso\textasciicircum{}r, put epsilon\_X := epsilon\_r/C\_rect; then 0 <= epsilon\_X <= e\_S1 and epsilon\_r = C\_rect*epsilon\_X. Consequently the scalar conclusions (R) of lem-stage1-polar-constant-ledger, with delta = delta\_* and r = rho, give every guard required in (A\_7): rho <= delta\_*, C\_ch*(epsilon\_r+delta\_*) <= kappa\_ch, C\_pol*(epsilon\_r+delta\_*) <= kappa\_pol, C\_grp*epsilon\_r < delta\_*-C\_pol*(epsilon\_r*delta\_*+delta\_*\textasciicircum{}2), C\_der*(epsilon\_r+rho) <= kappa\_der, and (1+epsilon\_r)*(1+C\_ch*(epsilon\_r+delta\_*))*rho+C\_grp*epsilon\_r < 2*delta\_*; moreover c := C\_der*(epsilon\_r+rho) <= kappa\_der/4 < 1. \steptag{validated} \\[4pt]
\step{1.2} For any algebra and epsilon\_r fixed as in 1.1, the objects used below define one smooth map sigma on all of calU. Indeed (A\_2) of lem-stage1-polar-constant-ledger supplies the unique C\textasciicircum{}1 graph family and gives ||D\_\{A\textasciicircum{}perp\}f\_V-I||<1, hence each D\_\{A\textasciicircum{}perp\}f\_V is invertible; lem-stage1-smooth-unitary-atlas upgrades exactly these graphs and charts to a smooth embedded atlas. Item (A\_4) supplies the C\textasciicircum{}1 polar diffeomorphism, and lem-stage1-smooth-polar-inverse upgrades that same map and same inverse to smooth maps. The guards from 1.1 make (A\_5) applicable, so its sigma(U)=u\_\{delta\_*\}(U\textasciicircum{}dagger) is global and sigma(J)=J. Applying lem-stage1-explicit-smooth-unitary-operations to these displayed graphs, atlas, polar map and inverse makes this same sigma smooth without changing its points or first derivatives, and its scalar covariance gives sigma(-J)=sigma((-1)J)=(-1)sigma(J)=-J. \steptag{validated} \\[4pt]
\step{1.3} For each s in \{+1,-1\}, let F\_s be the same-chart map supplied by (A\_7) of lem-stage1-polar-constant-ledger at delta=delta\_* and r=rho, and put G\_s:=F\_s-id on B\_rho\textasciicircum{}\{i calH\}(0). By 1.1, (A\_7) applies and gives chart retention and ||DG\_s(A)+2I|| <= c < 1 throughout that ball; by 1.2, G\_s is C\textasciicircum{}1. With V:=-2I, a Banach-space isomorphism of i calH, one has ||V\textasciicircum{}(-1)DG\_s(A)-I||=(1/2)||DG\_s(A)+2I|| <= c/2 < 1. Therefore lem-stage1-quantitative-inverse-function makes G\_s injective on B\_rho\textasciicircum{}\{i calH\}(0). Also f\_\{sJ\}(0)=0 using only exact unitality, bilinearity and s\textasciicircum{}2=1, so uniqueness in (A\_2) gives g\_\{sJ\}(0)=0 and chi\_s(0)=sJ; the fixed-center identities in 1.2 then give F\_s(0)=0 and G\_s(0)=0. Thus A=0 is the unique zero of G\_s in B\_rho\textasciicircum{}\{i calH\}(0). \steptag{validated} \\[4pt]
\step{1.4} Let U in calU satisfy sigma(U)=U and ||U-J||<rho. Exact unitality gives phi\_J(U)=U-J=:B; its anti-Hermitian and Hermitian parts A=B\textasciicircum{}parallel and B\textasciicircum{}perp each have norm at most ||B||, because dagger is isometric, so A lies in B\_rho\textasciicircum{}\{i calH\}(0) and both parts lie in the domains from (A\_2). Since U=J bold-dot (J+B) and U\textasciicircum{}dagger bold-dot U=J, the defining expression gives f\_J(B)=0; uniqueness in (A\_2) therefore yields B\textasciicircum{}perp=g\_J(A), hence U=chi\_+(A). The fixed-point equality and same-chart retention imply F\_+(A)=phi\_J\textasciicircum{}parallel(U)=A, so G\_+(A)=0; 1.3 forces A=0, and uniqueness gives g\_J(0)=0, whence U=J. Conversely J is fixed by 1.2, so J is the only fixed point in its ambient rho-ball. \steptag{validated} \\[4pt]
\step{1.5} Let U in calU satisfy sigma(U)=U and ||U+J||<rho. Since L\_\{-J\}=-I by exact unitality and bilinearity, phi\_\{-J\}(U)=-(U+J)=:B; its anti-Hermitian and Hermitian parts A=B\textasciicircum{}parallel and B\textasciicircum{}perp each have norm at most ||B||, so A lies in B\_rho\textasciicircum{}\{i calH\}(0) and both parts lie in the domains from (A\_2). One has U=(-J) bold-dot (J+B); the involution axiom gives (J+B\textasciicircum{}dagger) bold-dot (-J)\textasciicircum{}dagger=U\textasciicircum{}dagger, and unitarity gives f\_\{-J\}(B)=0. Uniqueness in (A\_2) therefore yields B\textasciicircum{}perp=g\_\{-J\}(A), hence U=chi\_-(A). The fixed-point equality and same-chart retention imply F\_-(A)=phi\_\{-J\}\textasciicircum{}parallel(U)=A, so G\_-(A)=0; 1.3 forces A=0, and uniqueness gives g\_\{-J\}(0)=0, whence U=-J. Conversely -J is fixed by 1.2, so -J is the only fixed point in its ambient rho-ball. \steptag{validated} \\[6pt]
\step{1.5.1} Assume U in calU, sigma(U)=U, and ||U+J||<rho. By validated node 1.2, the global smooth sigma fixes -J and the (A\_7) negative chart is the same chart used by sigma; by validated node 1.3, G\_-:=F\_- - id has 0 as its unique zero on B\_rho\textasciicircum{}\{i calH\}(0). Put B:=phi\_\{-J\}(U)=-(U+J), using L\_\{-J\}=-I from exact unitality and bilinearity, and A:=B\textasciicircum{}parallel. Isometry of dagger gives ||A||,||B\textasciicircum{}perp||<=||B||<rho, hence these variables are in the (A\_2) graph domains. From U=(-J) bold-dot (J+B), the involution axiom gives U\textasciicircum{}dagger=(J+B\textasciicircum{}dagger) bold-dot (-J)\textasciicircum{}dagger; since U is unitary, f\_\{-J\}(B)=0. The uniqueness clause of (A\_2) gives B\textasciicircum{}perp=g\_\{-J\}(A), so U=chi\_-(A). Since sigma(U)=U and (A\_7) retains sigma(chi\_-(A)) in this chart, F\_-(A)=phi\_\{-J\}\textasciicircum{}parallel(U)=A; thus G\_-(A)=0. Validated node 1.3 yields A=0, and its established g\_\{-J\}(0)=0 gives B=0 and U=-J. Conversely validated node 1.2 gives sigma(-J)=-J. Therefore -J is the unique fixed point of sigma in the ambient rho-ball about -J. \steptag{validated} \\[4pt]
\end{proofbox}

\end{document}
